\documentclass[11pt]{article}
% Standalone so it compiles on its own; the body between %%%BEGIN and %%%END
% can also be \input into the main paper.
\usepackage[T1]{fontenc}
\usepackage{underscore}
\usepackage{amsmath,amssymb,amsthm}
\usepackage{booktabs}
\usepackage[margin=1in]{geometry}
\usepackage{xcolor}
\newtheorem{theorem}{Theorem}
\newtheorem{lemma}{Lemma}
\newtheorem{definition}{Definition}
\newtheorem{remark}{Remark}
\title{\textbf{Soundness of the $N$-Hop Abstraction}\\\large A refinement theorem transferring routing safety to the concrete protocol for arbitrary path length}
\author{}\date{}
\begin{document}\maketitle
%%%BEGIN nhop-soundness

\noindent
This section upgrades the central claim from ``we verified an idealized $N$-hop
model'' to a \emph{theorem}: safety proved on the abstract, unbounded-length
routing model \emph{transfers} to the concrete, full-cryptography protocol, for
every path length $N$. The argument decomposes the problem into a part that is
cryptographic but \emph{local} (discharged in the concrete model at bounded
size) and a part that is unbounded but \emph{structural} (discharged in the
abstract model for all $N$), and composes them by a pen-and-paper refinement.

\subsection*{Two models over a shared action alphabet}

\begin{definition}[Concrete model $\mathcal{C}$]
The full-cryptography, fixed-position theory \texttt{multihop.spthy}: HTLC
offers are signed messages sent over an \texttt{Out}/\texttt{In} network under a
Dolev--Yao adversary with key compromise (\texttt{Compromise\_Ltk}). Forwarding
is authenticated by verifying the previous hop's signature.
\end{definition}

\begin{definition}[Abstract model $\mathcal{A}$]
The linear-fact routing theory (\texttt{multihop\_nhop\_fees.spthy}): the HTLC is
handed hop-to-hop as an internal linear fact \texttt{Route(prev,me,ptr,y)}
(an \emph{idealised channel}) rather than a signed broadcast. A single generic
forward rule fires at every hop, so paths have arbitrary length~$N$.
\end{definition}

Both theories emit the same safety-relevant action facts
$\Sigma = \{\,\texttt{ForwardHTLC},\ \texttt{Redeemed},\ \texttt{Refunded}\,\}$
with identical arities; the intermediary-safety properties are stated over
$\Sigma$ only. (The concrete model additionally emits the position-indexed
\texttt{HTLCForwarded1/2}; the projection below collapses both to the abstract
generic \texttt{HTLCForwarded}, but no safety lemma depends on the index.)

\subsection*{The trap, and the local authentication lemma}

\begin{remark}[Why abstraction is not automatically sound]
Internalising the network as a trusted \texttt{Route} fact \emph{removes}
Dolev--Yao power: $\mathcal{A}$ has strictly fewer behaviours than $\mathcal{C}$.
Safety on $\mathcal{A}$ therefore does not, by itself, imply safety on
$\mathcal{C}$. The gap is exactly the adversary's ability to forge or replay a
forwarded HTLC; closing it requires a cryptographic authentication guarantee in
$\mathcal{C}$.
\end{remark}

\begin{lemma}[\textsc{Auth}, per-hop authentication; machine-checked in $\mathcal{C}$]
\label{lem:auth}
For every forwarding step, the incoming HTLC accepted by a node traces back to a
genuine offer/forward from the \emph{adjacent} previous hop, unless that hop's
long-term key is compromised. Formally (one hop of the path), for the first and
second forwarding positions:
\[
\begin{aligned}
&\texttt{HTLCForwarded1}(F_1,F_2,S,p,y,v_{\mathit{in}},v_{\mathit{out}})
 \Rightarrow
 (\exists\, p',t.\ \texttt{HTLCOffered}(S,F_1,p',y,v_{\mathit{in}})@t)
 \ \lor\ \texttt{LtkCompromised}(S),\\
&\texttt{HTLCForwarded2}(F_2,R,F_1,p,y,v_{\mathit{in}},v_{\mathit{out}})
 \Rightarrow
 (\exists\, S,p',v.\ \texttt{HTLCForwarded1}(F_1,F_2,S,p',y,v,v_{\mathit{in}})@t)
 \ \lor\ \texttt{LtkCompromised}(F_1).
\end{aligned}
\]
\end{lemma}

\begin{proof}[Status]
Machine-checked in \texttt{multihop.spthy} as
\texttt{Forward1\_Requires\_Offer} (266 steps) and
\texttt{Forward2\_Requires\_Forward1} (274 steps), together with their
compromise-free variants \texttt{*\_Honest}. Each quantifies over a single
adjacent pair $(\mathit{prev},\mathit{me})$ and never over the path or $N$; the
forward rules verify only the previous hop's signature and read only the two
adjacent channels. Hence \textsc{Auth} is \emph{local}.
\end{proof}

\begin{remark}[Locality is what makes \textsc{Auth} tractable]
\label{rem:local}
A \emph{generic} single signed-forward rule that fires at every hop makes the
\texttt{'htlc'} term self-referential (a forwarded offer is the output of
forwarding another offer). Tamarin's origin analysis then fails to terminate at
the \emph{sources} stage: we verified that on such a rule \emph{no} lemma
completes---not \textsc{Auth}, not its honest variant, not even a two-hop
reachability witness---at 300\,s, and a hand-written \texttt{[sources]} lemma
does not rescue it (420\,s). \textsc{Auth} is provable precisely because, taken
per hop, it is \emph{bounded}: it looks back one adjacent step and does not
unfold the path. This is the crux of the whole argument.
\end{remark}

\subsection*{Projection and the refinement theorem}

\begin{definition}[Projection $\pi$]
$\pi$ maps a concrete trace to an abstract one by (i) keeping every
$\Sigma$-action verbatim and (ii) deleting all network/crypto events
(\texttt{Out}, \texttt{In}, \texttt{KU}, signature checks). Position-indexed
\texttt{HTLCForwarded1/2} are relabelled to the generic \texttt{HTLCForwarded}.
\end{definition}

\begin{lemma}[Per-hop simulation]
\label{lem:sim}
Under \textsc{Auth}, every concrete forwarding step has a matching abstract
step: an honestly-authenticated incoming HTLC (Lemma~\ref{lem:auth}) is exactly
what licenses the abstract \texttt{Route} hand-off, so $\pi$ of a concrete step
is a legal abstract step. The adversarial case (a forged/replayed HTLC) is ruled
out by \textsc{Auth} except under key compromise, which both models expose
identically via \texttt{LtkCompromised}.
\end{lemma}

\begin{theorem}[Soundness of the abstraction]
\label{thm:main}
Let $\varphi$ be a safety property over $\Sigma$ proved on $\mathcal{A}$ for all
path lengths $N$. Then $\varphi$ holds on $\mathcal{C}$: for every concrete trace
$t$, $\pi(t)$ is an abstract trace (Lemmas~\ref{lem:auth},~\ref{lem:sim}, by
induction over the forwarding steps of $t$), so if $\pi(t)$ satisfies $\varphi$
then so does $t$. Equivalently, every concrete violation projects to an abstract
violation; contrapositively, abstract safety transfers. Hence the intermediary%
-safety and atomicity properties, verified structurally on $\mathcal{A}$ for
arbitrary $N$, hold for the concrete signed-message protocol at arbitrary $N$.
\end{theorem}

\begin{proof}[Structure]
By induction on the number of forwarding steps in $t$. Base case: zero forwards,
$\pi(t)$ trivially abstract. Step: extend $t$ by one concrete forward; by
Lemma~\ref{lem:sim} (using Lemma~\ref{lem:auth} to discharge the forgery case)
its projection extends $\pi(t)$ by one legal abstract forward. The
$\Sigma$-actions are preserved verbatim, so any $\Sigma$-violation in $t$ appears
in $\pi(t)$. $\square$
\end{proof}

\subsection*{What is machine-checked, and what is pen-and-paper}

\begin{center}
\small
\begin{tabular}{@{}lll@{}}
\toprule
\textbf{Ingredient} & \textbf{Where} & \textbf{Status} \\
\midrule
\textsc{Auth}, per-hop, local & $\mathcal{C}$ (\texttt{multihop.spthy}) & machine-checked (266 / 274 steps) \\
All-$N$ routing safety ($\varphi$) & $\mathcal{A}$ (\texttt{*\_nhop\_fees}) & machine-checked (28/29; 30/30; 551-step \texttt{Intermediary\_Never\_Loses}) \\
$\Sigma$-alignment / projection $\pi$ & both & by inspection ($\Sigma$ shared verbatim) \\
Simulation + induction (Thm.~\ref{thm:main}) & meta & pen-and-paper \\
Generic single-rule \textsc{Auth} & --- & \emph{boundary}: sources non-termination (Rem.~\ref{rem:local}) \\
\bottomrule
\end{tabular}
\end{center}

\noindent
The meta-theorem is pen-and-paper by necessity: Tamarin cannot quantify over all
$N$ or relate two theories (the same reason key-management meta-theorems are
argued externally). Its two hypotheses, however, are both machine-checked, and
its soundness rests on the \emph{locality} of \textsc{Auth}
(Remark~\ref{rem:local})---which must be, and has been, confirmed by inspection
of every forward rule: each references only its adjacent pair. If any forward
rule depended on global path structure the argument would break; none does.

%%%END nhop-soundness
\end{document}
